% ============================================================
% Section 6 — Integration with the BX3 Framework
% ============================================================
\section{Integration with the BX3 Framework}
\label{sec:bx3-integration}

The Recursive Spawning Protocol is a structural instantiation of the \bxthree three-layer accountability architecture — not a module composed alongside it. Every RSP mechanism maps onto exactly one \bxthree layer, and each such mapping is structurally necessary: no layer in RSP is advisory, redundant, or replaceable by organizational convention. The immutable parent pointer, the depth bound, the escalation path, and the forensic ledger each operationalize a capability exclusive to a single layer. This section specifies the three mappings and demonstrates minimal-sufficiency.

% ------------------------------------------------------------
\subsection{Purpose Layer: Accountability Anchor}
\label{sec:purpose-anchor}

The Purpose Layer occupies two structurally exclusive roles in RSP, neither of which any other layer can perform.

\medskip
\noindent\textbf{Exclusive origin of spawn authorization.} At root provisioning, the Purpose Layer defines and records the spawn authorization policy $P_{\mathrm{spawn}}$ in the Fact Layer authorization ledger. $P_{\mathrm{spawn}}$ mandates two conditions for any valid child spawn: (i) $R(n_{i+1}) \subset R(n_i)$ — strict scope refinement — and (ii) operational necessity demonstrably unresolvable within the parent's scope. No machine actor may issue a spawn warrant in the absence of a matching Purpose Layer authorization record. The Bounds Engine evaluates against $P_{\mathrm{spawn}}$; the Fact Layer enforces it; neither originates it.

\medskip
\noindent\textbf{Exclusive terminus of escalation.} When a chain reaches $D_{\mathrm{ESCALATE}}$, the protocol compulsory-escalates to the human anchor $h(n)$. This terminus is non-collapsing: for all nodes $n_i \in C$, $h(n_i) = h(n_0)$. The human anchor does not migrate down the chain as depth increases. The human issues one of $\{\texttt{ACK}, \texttt{RESOLVE}, \texttt{REJECT}\}$; no machine actor may issue any of these. If any machine actor could absorb an exception or issue a terminal directive, the chain would terminate at that actor rather than at a human — defeating the upstream \bxthree accountability guarantee. The Purpose Layer's exclusivity here is not a design preference; it is the load-bearing constraint that makes the guarantee structurally operative.

% ------------------------------------------------------------
\subsection{Bounds Engine: Depth Enforcement}
\label{sec:bounds-depth}

The Bounds Engine enforces constrained execution evaluation: it evaluates proposed actions against the operating range defined by the Purpose Layer and enforces the depth constraint that prevents unbounded chain growth. The Bounds Engine cannot exceed its constraints under any operational circumstance.

\medskip
\noindent\textbf{Depth evaluation at spawn.} At each spawn event $n_i \xrightarrow{\text{spawn}} n_{i+1}$, the Bounds Engine evaluates current chain depth $|\text{depth}(n_i) + 1|$ against $D_{\max}$ (anchored in the Fact Layer authorization ledger). If $\text{depth}(n_i) + 1 > D_{\max}$, the spawn is refused, the Bounds Engine enters \texttt{ESCALATING} state, and the Bailout Protocol is triggered. This enforcement is deterministic: no discretion to exceed $D_{\max}$ regardless of urgency, criticality, or the absence of alternative resolution paths. The depth constraint is not a heuristic — it is a hard architectural boundary.

\medskip
\noindent\textbf{Fact Layer pre-commit reinforcement.} The Bounds Engine's depth evaluation is structurally reinforced by the Fact Layer pre-commit hook. Before any spawn is recorded, the Fact Layer re-verifies: (1) $R(n_{i+1}) \subseteq R(n_i)$ — scope refinement — and (2) $\text{depth}(n_{i+1}) \leq D_{\max}$ — depth constraint. Rejected spawns are logged; the Bounds Engine does not execute them; \texttt{ESCALATING} is triggered if the unmet condition persists. The pre-commit hook makes both the depth bound and scope subset relationship structural invariants — not policy conventions overridable in exigent circumstances.

The structural necessity of the Bounds Engine follows from the depth-limit insufficiency theorem (Section~\ref{sec:drift-depth-limits}): depth limits alone cannot prevent scope-creep drift when the operating scope is mutable. The Bounds Engine enforces the depth constraint; the Fact Layer enforces scope refinement; together they enforce both conjuncts of the drift prevention condition. Neither layer alone is sufficient.

% ------------------------------------------------------------
\subsection{Fact Layer: Forensic Ledger}
\label{sec:fact-ledger}

The Fact Layer provides deterministic enforcement and forensic auditability. It maintains the append-only ledger that makes the RSP chain auditable at any future time and enforces the write protocol that makes RSP tamper-evident.

\medskip
\noindent\textbf{Append-only forensic ledger.} The Fact Layer records every spawn event, every state transition, every exception condition, and every escalation action as an immutable ledger entry. Each entry contains: the acting node, the timestamp, the action type, the pre-state hash, and the post-state hash. Entries are hash-chained: each entry's integrity check field contains the cryptographic hash of the immediately prior entry, creating a tamper-evident sequence that propagates backward. Modifying any historical entry breaks the hash chain at the point of modification and is detectable by recomputing the chain integrity check from the genesis entry forward.

\medskip
\noindent\textbf{Immutable parent pointer establishment.} The Fact Layer performs the atomic write that establishes the parent pointer $\alpha(n_{i+1}) = n_i$ at spawn time. This write is immutable: no actor — including the Fact Layer itself after the write completes — can revise $\alpha(n_{i+1})$. The Fact Layer's deterministic execution rejects any subsequent write attempt to any established parent pointer field. This is not a privilege model; it is a structural property of the atomic write protocol.

\medskip
\noindent\textbf{Human anchor verification.} Before the first spawn event from any root node $n_0$, the Fact Layer verifies that $n_0$ has a designated human anchor $h(n_0)$ recorded in the authorization ledger. This verification is a pre-condition for chain initialization — not an advisory check. The chain terminates at $h(n_0)$, and no spawn event may occur in the absence of this anchor record. The Purpose Layer carries the human anchor identity as an immutable pre-loaded value; the Fact Layer carries the anchor record as an immutable ledger entry. Together they make the human anchor non-repudiable and是不可绕过的 (non-bypassable).

\medskip
\noindent\textbf{Structural necessity argument.} If the Fact Layer were mutable — if parent pointers could be revised post-spawn, or if spawn events could be recorded without hash-chained attestation — all three RSP correctness properties (Section~\ref{sec:rs-correctness}) collapse. Chain integrity depends on immutability; drift prevention depends on chain integrity; termination depends on the human anchor being reachable via an immutable chain. The Fact Layer is therefore not a logging convenience; it is the load-bearing enforcement substrate of the entire RSP correctness stack.

% ------------------------------------------------------------
\subsection{Minimal-Sufficiency of the Three-Layer Mapping}
\label{sec:layer-minimal-sufficiency}

Each RSP mechanism maps onto exactly one \bxthree layer, and no layer can substitute for another:

\begin{itemize}
    \itemsep0em
    \item \textbf{Purpose Layer} $\Rightarrow$ Spawn authorization origin and escalation terminus. No other layer can originate spawn policy or terminate escalation at a human.
    \item \textbf{Bounds Engine} $\Rightarrow$ Depth evaluation and scope constraint enforcement. The Purpose Layer defines scope; the Fact Layer records it; only the Bounds Engine evaluates scope compliance at spawn time and enforces depth as a hard boundary.
    \item \textbf{Fact Layer} $\Rightarrow$ Immutable parent pointer, hash-chained forensic ledger, human anchor verification. Neither Purpose Layer nor Bounds Engine can perform these operations without collapsing the immutability guarantee.
\end{itemize}

The three-layer RSP mapping is minimal in the sense that no layer is redundant, and sufficient in the sense that jointly they enforce all three correctness properties (Section~\ref{sec:rs-correctness}). This is not a design pattern — it is a structural necessity. Removing any layer causes the corresponding correctness property to fail.

% ============================================================
% Section 7 — Correctness Properties
% ============================================================
\section{Correctness Properties}
\label{sec:rs-correctness}

This section establishes three correctness properties for the Recursive Spawning Protocol — drift prevention, chain integrity, and termination — each stated formally, proved sketchedly, and connected to the \bxthree Framework layer that enforces it. Together these properties constitute the RSP's overall guarantee: a spawn chain is an accountable, bounded, and auditable refinement process that terminates in human judgment.

% ------------------------------------------------------------
\subsection{Drift Prevention}
\label{sec:rs-drift-prevention}

\medskip
\noindent\textbf{Theorem (Drift Prevention).} Let $C = \langle n_0, n_1, \ldots, n_k \rangle$ be a Recursive Spawning chain rooted at $n_0$ with human anchor $h(n_0) = h(n_k)$. Then for all $i$ ($0 \leq i \leq k$), the operating scope $R(n_i)$ is a descendant refinement of the intent of $h(n_0)$. Equivalently: no node in $C$ can exhibit behavior outside the authorized intent of its human anchor.

\medskip
\noindent\textbf{Proof sketch.} The proof proceeds by structural induction on chain depth, where each inductive step is enforced at runtime by the Fact Layer pre-commit hook — making drift prevention a structural invariant rather than an inferred property.

\medskip
\textit{Base case ($i=0$).} $R(n_0)$ is defined by $h(n_0)$ at provisioning. By the Purpose Layer's definition of authorized operating scope, $R(n_0)$ is a descendant refinement of $h(n_0)$'s intent by construction. No drift exists at the root.

\medskip
\textit{Inductive step.} Assume $R(n_i)$ is a descendant refinement of $h(n_0)$'s intent for some $i$ with $0 \leq i < k$. For $n_i$ to spawn $n_{i+1}$, the Fact Layer pre-commit hook must verify before recording:

\begin{enumerate}
    \item[(a)] $R(n_{i+1}) \subseteq R(n_i)$ — scope refinement. The Purpose Layer's spawn policy $P_{\mathrm{spawn}}$ requires strict scope subset; no authorized provision for lateral or upward scope expansion exists.
    \item[(b)] $\text{depth}(n_{i+1}) \leq D_{\max}$ — depth constraint, also enforced by the pre-commit hook.
\end{enumerate}

The Fact Layer rejects any spawn violating (a) or (b). No actor — including the Bounds Engine — can execute a rejected spawn. Therefore $R(n_{i+1}) \subseteq R(n_i)$ holds by structural enforcement.

By the inductive hypothesis, $R(n_i)$ refines $h(n_0)$'s intent. Since $R(n_{i+1}) \subseteq R(n_i)$, $R(n_{i+1})$ also refines $h(n_0)$'s intent. Hence $n_{i+1}$ does not drift.

\medskip
The proof sketch rests on three architectural conjuncts that collectively make drift structurally impossible:

\begin{enumerate}
    \item[(1)] \textbf{Immutable parent pointer.} $\alpha(n_{i+1}) = n_i$ is established atomically at spawn and can never be revised. The chain topology is permanent; no retroactive chain manipulation is possible.
    \item[(2)] \textbf{Forensic ledger with pre-commit enforcement.} The Fact Layer pre-commit hook verifies $R(n_{i+1}) \subseteq R(n_i)$ and $\text{depth}(n_{i+1}) \leq D_{\max}$ before any spawn is recorded. This enforcement is not a policy check — it is a deterministic atomic operation that fails (and blocks the spawn) if the condition is unmet.
    \item[(3)] \textbf{Chain terminates at Purpose Layer human anchor.} For all $n_i \in C$, $h(n_i) = h(n_0)$ is a human principal. No machine actor can serve as a terminal node in the accountability chain. If an exception cannot be resolved within $R(n_i)$, the Bailout Protocol compulsory-escalates to $h(n_i)$ — never to an intermediate machine actor.
\end{enumerate}

If any one of these conjuncts fails, drift becomes possible. If the parent pointer were mutable, an adversarial actor could redirect the chain. If the forensic ledger were mutable or absent, scope refinement could not be enforced. If the chain could terminate at a machine actor, a machine actor could absorb exceptions and continue operating without human authorization. The conjunction of all three makes drift architecturally impossible. $\square$

% ------------------------------------------------------------
\subsection{Chain Integrity}
\label{sec:rs-chain-integrity}

\medskip
\noindent\textbf{Theorem (Chain Integrity).} For any Recursive Spawning chain $C = \langle n_0, n_1, \ldots, n_k \rangle$, the parent pointer sequence $\alpha(n_i) = n_{i-1}$ holds for all $1 \leq i \leq k$, and no node configuration in $C$ has been modified after provisioning.

\medskip
\noindent\textbf{Proof sketch.} Chain integrity follows from two independent enforcement mechanisms:

\medskip
\textit{Immutability of the parent pointer.} The Fact Layer establishes $\alpha(n_i) = n_{i-1}$ via an atomic write at spawn time. The Fact Layer's deterministic execution rejects any subsequent write to an established parent pointer field. Since the Fact Layer is the sole write authority for parent pointer fields, and its write protocol is immutable post-establishment, $\alpha(n_i) = n_{i-1}$ holds for all $i$ by construction. No runtime operation — including exceptions, escalations, or system faults — can modify $\alpha$.

\medskip
\textit{Hash-chained forensic ledger.} Every node provisioning event, state transition, and chain event is recorded in the Fact Layer forensic ledger with a hash-chained integrity field. Each entry $e_j$ contains $\text{hash}(e_{j-1})$ as its integrity check value. Any modification to a historical entry breaks the hash chain at the point of modification and is detectable by recomputing the chain forward from the genesis entry. The ledger records all topology-affecting events; any unauthorized modification is immediately detectable by any auditor.

\medskip
\textit{Node configuration immutability post-provisioning.} The configuration of any node $n_i$ — its operating scope $R(n_i)$, its spawn parameters, its exception handling policy — is sealed at provisioning and cannot be modified by any machine actor post-establishment. Modifications require a new spawn event with a new Fact Layer entry, creating a new node with a new configuration. The original node's configuration is immutable and auditable.

The conjunction of these three mechanisms ensures that the chain topology and all node configurations are stable from provisioning forward. Chain integrity is therefore a structural property, not a policy assertion. $\square$

% ------------------------------------------------------------
\subsection{Termination}
\label{sec:rs-termination}

\medskip
\noindent\textbf{Theorem (Termination).} Every Recursive Spawning chain $C = \langle n_0, n_1, \ldots, n_k \rangle$ terminates — either at a resolution state or at the escalation threshold — within a finite number of spawn steps bounded by $\max(D_{\max}, D_{\mathrm{ESCALATE}})$.

\medskip
\noindent\textbf{Proof sketch.} Two lemmas combine to establish the theorem.

\medskip
\textit{Lemma 1 (Depth-bounded growth).} The chain grows only through spawn events. By the Bounds Engine enforcement and the Fact Layer pre-commit hook, any spawn event producing $\text{depth}(n_i) + 1 > D_{\max}$ is rejected. Autonomous chain growth is thus bounded by $D_{\max} - 1$ spawn steps. Rejected spawns are logged; the Bounds Engine transitions to \texttt{ESCALATING}.

\medskip
\textit{Lemma 2 (Escalation-triggered halt).} If the chain reaches $D_{\mathrm{ESCALATE}}$ before resolving the exception condition, the protocol enters \texttt{ESCALATING} and suspends all autonomous spawn activity. The Bounds Engine awaits a directive from $h(n_k)$ — the human anchor. Since $h(n_k) = h(n_0)$ is a human principal, the directive arrives in finite time. \texttt{REJECT} terminates the chain definitively; \texttt{RESOLVE} redirects behavior and records a new ledger entry; \texttt{ACK} permits continued operation under continued Purpose Layer authorization.

\medskip
\textit{Theorem conclusion.} The chain grows autonomously at most $D_{\max} - 1$ spawn steps. If it reaches $D_{\mathrm{ESCALATE}}$ before resolution, it halts and escalates. If it reaches $D_{\max}$ before resolution, the Fact Layer rejects further spawn and triggers mandatory escalation. In all execution paths, the chain reaches a terminal state within $\max(D_{\max}, D_{\mathrm{ESCALATE}})$ spawn steps. Since $D_{\max}$ and $D_{\mathrm{ESCALATE}}$ are finite integers defined at provisioning, the chain terminates in finite time. $\square$

The termination guarantee also precludes infinite spawn-escalate cycles. When $D_{\max}$ is reached, the Fact Layer blocks further spawn and triggers mandatory escalation. The human anchor's \texttt{REJECT} directive terminates the chain. No mechanism exists by which the chain can circumvent this bound and re-enter autonomous spawn — the Fact Layer pre-commit hook is structural, not configurable at runtime by any machine actor.

% ------------------------------------------------------------
\subsection{Collective Guarantee}
\label{sec:rs-collective}

Together, the three correctness properties constitute RSP's overall accountability guarantee in the context of recursive agent population dynamics.

Drift prevention ensures scope refinement is the only authorized direction of chain growth. Without it, successive scope expansions could gradually migrate the chain's operating scope away from the human anchor's intent — an autonomy runway within a bounded-depth chain. The depth-limit insufficiency theorem (Section~\ref{sec:drift-depth-limits}) establishes that depth bounds alone cannot prevent this.

Chain integrity ensures chain topology is stable and tamper-evident. The parent pointer sequence is immutable, node configurations are unmodifiable post-provisioning, and the forensic ledger's hash-chaining makes any tampering detectable. Without chain integrity, an adversarial actor could manipulate chain parentage or configuration to redirect accountability or obscure node provenance — a form of accountability laundering.

Termination ensures the chain cannot grow indefinitely and cannot sustain an infinite spawn-escalate cycle. The depth bound is enforced structurally by the Fact Layer pre-commit hook, and the escalation path terminates in human judgment, not in autonomous resolution. Without termination, an unresolvable exception condition could drive unbounded chain growth exceeding any human's supervisory capacity — defeating the upstream \bxthree accountability guarantee by overwhelming the human anchor.

These three properties are mutually reinforcing and jointly enforced by the \bxthree three-layer architecture. The Purpose Layer defines spawn authorization policy and receives escalated exceptions — the exclusive terminus of every accountability path. The Bounds Engine evaluates spawn necessity and enforces the depth constraint — it cannot bypass the depth bound through operational urgency. The Fact Layer enforces the immutable parent pointer, maintains the append-only forensic ledger, and executes the pre-commit hooks that make all other constraints structural invariants. This mapping is not a design pattern — it is a structural necessity. Removing any layer causes the corresponding guarantee to collapse.

The result is an accountable, bounded, and auditable refinement process: it grows only in authorized directions (drift prevention), preserves the topology established at provisioning (chain integrity), halts within a Purpose Layer-defined depth bound (termination), and maintains a complete forensic record of every decision point (Fact Layer). This is the operational expression of the upstream \bxthree accountability guarantee applied to recursive agent population dynamics.